% Om Shri Ganeshaye Namah, Om Shri Saraswattye Namah
%\documentclass[showpacs,aps,amssymb,floatfix,prd,amsmath,twocolumn,preprintnumbers]{revtex4}
\documentclass[showpacs,aps,amssymb,preprintnumbers,prd,amsmath,floatfix]{revtex4}
\setlength{\textheight}{10in} % Used for PRD
\usepackage{graphicx}  
\usepackage{dcolumn}   
\usepackage{bm}
\usepackage{color}
\usepackage{amsmath}
\usepackage{amssymb}
\usepackage{epsfig}
\usepackage{amsfonts}
\usepackage{lineno,hyperref}
\usepackage{array}
\usepackage{float}
\usepackage{microtype}
\usepackage{multirow}
\usepackage{adjustbox}
\usepackage{blindtext}
\usepackage{caption}
\usepackage{subcaption}
\usepackage{orcidlink}



\def \m{\mu}
\def \n{\nu}
\def \a{\alpha}
\def \b{\beta}
\def \g{\gamma}
\def \d{\delta}
\def \l{\lambda}
\def \e{\epsilon}
\def \k{\kappa}
\def \s{\sqrt}
\def \be{\begin{equation}}
\def \ee{\end{equation}}
\def \ben{\begin{eqnarray}}
\def \een{\end{eqnarray}}
\def \o{\omega}
\def \nb{\nabla}
\def \sg{\sigma}
\def \p{\partial}
\def \t{\theta}
\def \M{\mathcal{M}}
\def \L{\mathcal{L}_m}
\def \A{\mathcal{A}}
\def \r{\rho}
\def \G{\Gamma}
\def \non{\nonumber}
\def \T{\Theta}
\def \red{\textcolor{red}}
\def \rf{\tilde{\r}}
\def \pf{\tilde{p}}
\def \sh{\sech}
\def \ch{\cosh}
\def \th{\tanh}
\def \blue{\textcolor{blue}}



%\usepackage{nicefrac} 
%\usepackage{subcaption}
%%%%%%%%%%%%%%%%%%%%%%%%%%%%%
\begin{document}
	
%\input epsf.tex
%%%%%%%%%%%%
%%%%%%%%%%%

\title{Aspects of Cosmology in Symmetric Teleparallel $f(Q,T)$ Gravity}

\author{Surajit Das\orcidlink{0000-0003-2994-6951}$^{1}$\footnote{Email: surajit.cbpbu20@gmail.com},
  Sanjay Mandal\orcidlink{0000-0003-2570-2335}$^{2}$\footnote{Email: sanjaymandal960@gmail.com
  ~(corresponding author)}}\
  
\affiliation{$^{1}$ Department of Physics, Cooch Behar Panchanan Barma University, Panchanan Nagar, Vivekananda Street, Coochbehar, West Bengal-736101, India,\\
$^{2}$ Faculty of Mathematics and Computer Science, Transilvania University, Iuliu Maniu Str. 50, 500091 Brasov, Romania.}



\begin{abstract}
This article aims to explore the idea of describing the complete evolution process of the universe with a cosmological model. For this purpose, we work on a recently developed generalized symmetric teleparallel gravity called $f(Q,T)$ gravity, where $Q$ and $T$ represent the non-metricity scalar and trace of the energy-momentum tensor. We present two $f(Q,T)$ cosmological models for two different Lagrangian forms of $f(Q,T)$ and discuss their cosmological parameters. Further, we discuss various energy conditions to check the viability of those models and also discuss the scalar field construction for describing the inflationary scenario in the early times. In addition, we examine the dark energy nature of the presumed cosmological solution by employing state-finder diagnostics and $\o-\o^{\prime}$ phase space analysis. In conclusion, we found that our models can present a complete evolution profile from the early to late times of the universe.
\\ \\
\textbf{Keywords:} $f(Q,T)$ gravity, cosmic kinematic parameters, energy conditions (ECs), scalar field model, statefinder diagnostics.
\end{abstract}


\pacs{04.50.Kd}



\maketitle






\section{Introduction}
Due to glorious and advanced progress in astronomy, astrophysics, cosmology, data science, space science \& technology, recent observations suggest that our universe is accelerating with unknown stuff, so-called `Dark Energy', which is responsible for the acceleration of the universe \cite{Riess1,Perlmutter1,Perlmutter2,Spergel1,Komatsu}. But the most intriguing question is that what is dark energy? Cosmologists and particle physicists are trying to identify the properties of dark energy. Our universe is covered with approximately 75\% of the total unknown energy i.e., dark energy. So it's a big challenge for the physicists in this field to overcome the related problems of dark energy. There exist two eras in the cosmic history of our universe for accelerated expansion \cite{Capoz}. During the early times, the universe passed through an accelerated expansion, which is known as inflation. And now, in the current times, the expansion of the universe is accelerating due to the dominating dark energy.

Historically, it was added a term, called the cosmological constant ($\Lambda$) in Einstein's field equation described correctly at the very first time, but its magnitude is unwanted and unmotivated by the law of fundamental physics. One may explain both, the inflationary epoch and the late-time accelerated expansion of the universe by considering matter which contains barotropic fluid having the form $p=f(\r)$, where $p$ and $\r$ are the pressure and the energy density respectively. The study of cosmological regimes beyond the phantom barrier and future singularity has been discussed in \cite{Nojiri1,Nojiri2} as an example of barotropic fluids. Also, the inflationary scenarios in the viscous fluid model have been described in \cite{Bamba}. On the other hand, in this literature \cite{Astas}, they argued that one cannot sure to determine the ultimate fate of our universe when we are using the current observational data. So looking at all these problems, researchers have proposed a lot of proposals to overcome the problems related to dark energy \cite{Ratra,Caldwell1,Buchert,Picon,Tomita,Milton,Hunt,Easson,Radicella,Pavn,Pandey}. 

The modification of the standard general theory of relativity is a great idea to describe dark energy. Therefore the modified theories of gravity are interesting proposals to study the existence of dark energy. It represents the nature of dark energy as a geometrical property of the universe and the modified gravity theories are connected to the modification of Einstein-Hilbert action \cite{Palia}. There are a lot of interesting modified theories of gravity that exist. For example, in the Einstein-Hilbert action, if we replace by an analytic function $f(R)$ instead of $R$, then we get $f(R)$ gravity \cite{Soti}. On the other hand, in between geometry-matter coupling in a non-minimal way, we have another interesting modified theory of gravity, known as $f(R,T)$ gravity \cite{Har}. Another well-established modified theory of gravity is $f(T)$ gravity, which is also known as teleparallel equivalent general relativity (TEGR) \cite{Bengochea,Ferraro,Linder1,Cai}. In teleparallel gravity the action is an arbitrary function of torsion scalar $T$, where Ricci scalar $R$ is replaced by the arbitrary function of torsion scalar $T$ in the action.

Besides the two representations in general relativity (one is curvature representation and another is torsion representation), another interesting and well-established representation in the modified theory of gravity, which has attracted the interests of cosmologists, is the so-called symmetric teleparallel gravity \cite{Nester}. In this formalism, the geometric variable is defined by a non-metricity $Q$ and in this aspect, $f(Q)$ gravity has been proposed \cite{Lavinia}. Like $f(R,T)$ gravity, one may extend the $f(Q)$ theory by considering the action which is a non-minimal coupling in between gravitational interaction produced by the term $Q$ and the trace of the energy-momentum tensor $T$. So it is an action containing an arbitrary function of $f(Q,T)$. In $f(Q,T)$ gravity, the study of cosmological scenarios \cite{Nisha,Xu,Gad}, constraining equation of state \cite{Simran3}, cosmological perturbation \cite{Antonio}, inflationary model \cite{Maryam}, Quintessence universe \cite{Koussour}, late-time cosmic acceleration \cite{Lalke}, energy conditions \cite{Simran1,Simran2}, wormholes \cite{Tayde}, holographic dark energy \cite{Asutosh}, reconstruction of the $\Lambda$CDM universe \cite{gaurav} has been studied.

In the present paper, we proceed in a phenomenological way to study the complete cosmic evolutionary history of the universe by describing an ansatz for the scale factor. From this ansatz, one can easily find out the analytical solutions for energy density, pressure, EOS parameter, and various interesting physical consequences for various cosmological models in $f(Q,T)$ gravity. To explain the early as well as late-times cosmic acceleration of the universe from the chosen parametrization, we have studied all the cosmological consequences in detail in the framework of $f(Q,T)$ gravity. 

This paper is presented as follows: We started with a discussion on the STEGR in section \ref{II}. In section \ref{III}, we have discussed in short and briefly the mathematical formulation of $f(Q,T)$ gravity and presented the gravitational field equations in a spatially flat FLRW-type spacetime. In section \ref{IV}, we discuss some kinematic cosmic variables such as scale factor, Hubble parameter, and deceleration parameter, which is followed by the construction of the cosmological model in $f(Q,T)$ gravity in section \ref{V} for two different models. Here we present the expressions and behaviors for energy density, pressure, and equation of state parameters for the given models. The energy conditions for the chosen models are presented in section \ref{VI}. The formulation of the scalar field description for our models is discussed in section \ref{VII}. In sections \ref{VIII} \& \ref{IX}, we have discussed some geometrical diagnostics to distinguish our model from other dark energy models and finally, we conclude our results in \ref{X} with a brief discussion.






\section{Symmetric Teleparallel Equivalent General Relativity in a Nutshell}\label{II}
Since general relativity (GR) is based on the Riemannian manifold, it should be generalized to a more general geometric theory of gravity which could explain the gravitational field as a more general geometric gravitational structure that should be valid in the solar system level also. So we may explain the various cosmological aspects of the universe from the early to the late time acceleration. In this generalization process, Weyl introduced the notion of non-metricity, a new geometric quantity, in which the covariant derivative of the metric tensor is non-zero \cite{Weyl/1918}. In Weyl's theory, there is an extra connection, called the length connection, which gives us information on the length of a vector and doesn't contain any knowledge about the direction of a parallel transported vector field. After Weyl, more generalization of gravity theories has been formulated and can be found here \cite{Di,Von,D,Z,M1,T,L,J,Z1,Z2}.

So from the previous discussions, it is clear that GR can be presented at least in two equivalent formalisms. One is the well-known curvature representation, where the non-metricity and the torsion tensor vanish. Another one is called the teleparallel equivalent general relativity (TEGR) in which the curvature and the non-metricity vanish, but the torsion tensor does not. But fortunately, a relatively unexplored territory consists of another third equivalent representation of GR, namely symmetric teleparallel equivalent of general relativity (STEGR) or simply symmetric teleparallel gravity \cite{Nester}. In this theory, the curvature and the torsion vanish but the non-metricity ($Q$) is the basic geometrical variable that describes the properties of geometric gravitational interactions. It also describes the length-variation of a parallel transported vector. This formulation is completely geometric and covariant. As the curvature and torsion relate with affine connections, not dealing with spacetime, the covariant derivatives must commute due to the vanishing of curvature. Here in STEGR, the associated energy-momentum density is essentially the Einstein pseudotensor but becomes a true tensor in the geometric representation. STEGR was further developed into an arbitrary generic function of $Q$ i.e., $f(Q)$ gravity which is also known as coincident general relativity (CGR) \cite{Lavinia}. CGR is described by the Einstein-Hilbert action excluding the boundary term which is underpinned by the spin-2 field theory. This particular construction also provides a starting
point for the modified gravity theories, and presents the early and late-times cosmological solutions of the universe. That's why this is a simpler geometrical structure to the affine connection which is fundamentally devastating gravity from any inertial character.

The STEGR can also be represented by a general quadratic and parity-conserving Lagrangian with the help of Lagrangian undetermined multipliers for vanishing torsion and curvature \cite{Adak}. This particular Lagrangian is equivalent to the Einstein-Hilbert Lagrangian in standard GR for certain choices of the coupling coefficients. However, it was also shown that the field equations can be written as a system of Proca equations, which may be an interesting study for the propagation of gravitational-electromagnetic waves \cite{Mol}. Conroy et. al., \cite{Conroy} studied the action, completely made up of the non-metricity tensor, and its contractions were decomposed into terms, which involved the metric and a gauge vector field. So this paper discussed the derivation of the exact propagator for the most general infinite-derivative, even-parity in the generally covariant STEGR theory of gravity. Also, the linear perturbations in flat space were analyzed in \cite{Conroy,Kovisto}, and \cite{Beltran}. The propagation of gravitational waves with their various properties like speed and polarization was studied in \cite{Soudi} for the various extensions of symmetric teleparallel gravity. For all the possible classifications of quadratic, first-order derivative terms of the non-metricity tensor in the framework of symmetric teleparallel gravity, Dialektopoulos et. al.,\cite{Dia} used the Noether symmetry approach. Basically in \cite{Dia}, they used to reduce the dynamics of the system by choosing symmetries to find analytical solutions and this model were invariant under point transformations in a cosmological background. 

On the other hand, from the cosmological point of view, it was seen in the pieces of literature \cite{Lu} and \cite{Laz} that, the accelerating expansion of the universe is an intrinsic property of geometry and we need not deal with any extra fields or exotic dark energy, under $f(Q)$ gravity consideration. Cosmology and the behavior of cosmological perturbations in $f(Q)$ gravity were investigated in \cite{Lavi2}. Energy conditions, cosmography analysis, Buchdahl quark star formation, and wormhole solutions in $f(Q)$ gravity can be found in \cite{Sanjay1}, \cite{Sanjay2}, \cite{Sneha}, and \cite{Mustafa1,Hassan1} respectively. Some excellent works in the aspects of cosmology can be found here \cite{Raja,De,Raja2,Sanjay3,Sanjay4}.

By introducing in the framework of the metric-affine formalism, an extension of a new class of symmetric teleparallel gravity was considered in \cite{Harko}, where the geometry part i.e., the non-metricity $Q$ is non-minimally coupled to the matter Lagrangian $L_{m}$. A Lagrangian of the form $L=f_1(Q)+f_2(Q)L_{m}$ leads to the non-conservation of the energy-momentum tensor and appearance as an extra force in the geodesicity. Here $f_1$ and $f_2$ are two generic functions of $Q$, and $L_{m}$ is the considered matter Lagrangian. Several cosmological applications were considered for some specific functional forms of the functions like power-law and exponential dependencies of the non-minimal couplings and the cosmological solutions lead to a late-times accelerating universe.

Finally, one may consider another kind of extension of symmetric teleparallel gravity theory by considering an action that contains the non-minimal coupling between the geometry i.e., $Q$, and the trace of the energy-momentum tensor $T$ instead of matter Lagrangian unlike in the previous one. By considering this kind of construction, one can easily describe the gravitational interactions in the presence of geometry-matter coupling and the cosmological solutions can describe both, the accelerating and decelerating evolutionary phases of the universe. So $f(Q,T)$ gravity can provide hopeful insights for the description of the early and late phases of our universe.






\section{Mathematical Formulation of $f(Q,T)$ Gravity}\label{III}
In the symmetric teleparallel $f(Q,T)$ gravity theory, by imposing the condition that the connection is symmetric, so that for coincident gauge, the Levi-Civita connection $\hat{\G}^\l_{\m\n}=0$. So the disformation tensor $L^\l_{\m\n}$ can be written as,
\be
L^\l_{\m\n}=-\G^\l_{\m\n}\label{1},
\ee
The non-metricity tensor is given by $Q_{\a\m\n}=\nb_\a g_{\m\n}$ and it can be also defined as,
\be
Q=-g^{\m\n}\Big(L^\a_{\b\m}L^\b_{\n\a}-L^\a_{\b\a}L^\b_{\m\n}\Big)\label{2},
\ee
where from \eqref{1}, the disformation tensor is written as,
\be
L^\a_{\b\g}=-\frac{1}{2}g^{\a\l}\Big(-\nb_\l g_{\b\g}+\nb_\b g_{\g\l}+\nb_\g g_{\b\l}\Big).\label{3}
\ee
The non-metricity conjugate or the superpotential of the model is given by,
\be
P^\a_{\m\n}=-\frac{1}{2}L^\a_{\m\n}+\frac{1}{4}(Q^\a-\Bar{Q}^\a)g_{\m\n}-\frac{1}{4}\d^\a~_{(\m}Q_{\n)},\label{4}
\ee
where we have used two known results of the trace of the non-metricity tensor i.e.,
\be
Q_\a=g^{\m\n}Q_{\a\m\n}\equiv Q^{~\m~}_{\a~~\m}~;~~~~
\bar{Q}_\b=g^{\m\n}Q_{\m\b\n}\equiv Q^{\m}_{~~\b\m}.\label{5}
\ee
In addition, we define the non-metricity scalar $Q$ as,
\be
Q=-Q_{\a\m\n}P^{\a\m\n}=-\frac{1}{4}\Big(-Q^{\a\n\l} Q_{\a\n\l}+2 Q^{\a\n\l}Q_{\l\a\n}
-2 Q^\l\bar{Q}_\l+Q^\l Q_\l\Big).\label{6}  
\ee
Now the corresponding action in $f(Q,T)$ gravity can be written as,
\be
\A=\int \frac{1}{16\pi G}f(Q,T)\s{-g}~d^4x+\int \L\s{-g}~d^4x\label{7},
\ee
where $f$ is a generic arbitrary function of $Q$ and the trace of the energy-momentum tensor is $T$, $G$ being the Newtonian universal gravitational constant. The matter Lagrangian can be considered as $\L$ and $g$ is the determinant of the metric tensor $g_{\m\n}$. Here the energy-momentum tensor is defined as,
\ben
T_{\a\b}=-\frac{2}{\s{-g}}\dfrac{\d({\s{-g}\L})}{\d g^{\a\b}}.\label{8}
\een
So taking the variation of the action \eqref{7} wrt the metric tensor $g_{\m\n}$, we obtain the field equation of the $f(Q,T)$ gravity as
\ben
-\frac{2}{\s{-g}}\nb_\l\Big(f_Q\s{-g}P^\l_{~~\m\n}\Big)-\frac{1}{2}f g_{\m\n}-f_Q\Big(P_{\m\l\b}~Q^{~~\l\b}_\n
-2 Q^{\l\b}_{~~\m}P_{\l\b\n}\Big)=\big(8\pi G-f_T\big)T_{\m\n}-f_T\T_{\m\n},\label{9}
\een
where $f_Q=\frac{df}{dQ}$ and $f_T=\frac{df}{dT}$. $\T_{\m\n}$ is the variation of energy-momentum tensor wrt the metric tensor, such that
\ben
\T_{\m\n}\equiv g^{\a\b}\frac{\d T^{\a\b}}{\d g^{\m\n}}.\label{10}
\een
It should be mentioned that the field equation, \eqref{9} is valid only for the coincident gauge \cite{Lavinia}.

By assuming the universe is homogeneous and isotropic, we consider flat FLRW-type metric having the form
\ben
ds^2=-dt^2+a^2(t)\sum_i(dx^{i})^2.\label{11}
\een
Here $a(t)$ describes the evolution history of the universe, simply known as scale factor and $i$ runs from $1-3$ only. Moreover, considering the cosmic fluid behaves as perfect fluid having the energy-momentum tensor as,
\ben
T_{\m\n}=(\r+p)u_\m u_\n +pg_{\m\n},\label{12}
\een
where $\r$, $p$ and $u^\m$ are the energy density, pressure and the four-velocity of the cosmic fluid, respectively. The non-metricity $Q$ in this background can be calculated and it is given by $Q=6H^2$.

Using the metric \eqref{11} in the field equations \eqref{9}, we can write the generalized Friedmann equations (taking $G=1$), which are
\ben
8\pi \r=\frac{f(Q,T)}{2}-6FH^2-\frac{2\bar{G}}{1+\bar{G}}\big(\dot{F}H+F\dot{H}\big),\label{13}
\een
and
\ben
8\pi p=-\frac{f(Q,T)}{2}+6FH^2+2\big(\dot{F}H+F\dot{H}\big)\label{14}.
\een
Here (·) dot represents a derivative with respect to cosmic time, $F=\frac{df}{dQ}$ and $8\pi\bar{G}=\frac{df}{dT}$.
\\
We can also write Einstein’s field equations from Friedmann equations as follows
\begin{equation}
3H^2\equiv 8\pi \r_{eff}=\frac{f}{4F}-\frac{4\pi}{F}\big[(1+\bar{G})\r+\bar{G}p\big],\label{15}
\end{equation}
\ben
2\dot{H}+3H^2\equiv -8\pi p_{eff}=\frac{f}{4F}-\frac{2\dot{F}H}{F}
+\frac{4\pi}{F}\big[(1+\bar{G})\r+(2+\bar{G})p\big]\label{16}.
\een
Here $\r_{eff}$ and $p_{eff}$ are the effective energy density and effective pressure, respectively.
\\
Using \eqref{13} and \eqref{14}, the equation of state parameter can be written as,
\ben
\o=-1+\dfrac{4(\dot{F}H+\dot{H}F)}{f(1+\bar{G})-12FH^2(1+\bar{G})-4\bar{G}(\dot{F}H+\dot{H}F)}.\label{17}
\een

Now using the above setup, one can explore various cosmological scenarios in the framework of $f(Q,T)$ gravity.






\section{Cosmic Parameters}\label{IV}
In this section, we shall discuss about various cosmological parameters like the scale factor $a(t)$, the Hubble parameter $H(t)$ and the deceleration parameter $q(t)$. These particular kinematic variables play very crucial roles in the study of physical cosmology. Also these parameters are the key parameters of most of the various cosmological models in modified gravity theories and these can describe the evolutionary history of the universe. 

In this paper, we choose the scale factor $a(t)$ having the following form \cite{Moraes}, 
\begin{equation}
a(t)=e^{ct}[\text{sech}^d(n-mt)].\label{18}
\end{equation}

\begin{figure}[H]
  \centering
  \includegraphics[width=8.5 cm]{avst.eps}
  \caption{Plot of scale factor $(a)$ as a function of cosmic time $(t)$ for $c=0.97,d=1,m=0.735,$ \& $n=10$.}
  \label{f1}
\end{figure}

Using \eqref{18}, one can obtain the Hubble parameter $H(t)$ as
\begin{equation}
H(t)=\frac{\dot{a}}{a}=c+dm \tanh(n-mt).\label{19}
\end{equation}

\begin{figure}[H]
  \centering
  \includegraphics[width=8.5 cm]{Hvst.eps}
  \caption{Plot of Hubble parameter $(H)$ as a function of cosmic time $(t)$ for $c=0.97,d=1,m=0.735,$ \& $n=10$.}
  \label{f2}
\end{figure}

Also the deceleration parameter $q(t)$ is given by,
\begin{equation}
q=-1-\frac{\dot{H}}{H^2}=-1+\frac{d m^2}{[c\cosh(n-mt)+dm\sinh(n-mt)]^2}\label{20}.
\end{equation}

\begin{figure}[H]
  \centering
  \includegraphics[width=8.5 cm]{qvst.eps}
  \caption{Plot of deceleration parameter $(q)$ as a function of cosmic time $(t)$ for $c=0.97,d=1,m=0.735,$ \& $n=10$.}
  \label{f3}
\end{figure}

It is worthy mentioned from the analysis of the research articles \cite{Riess1,Perlmutter1,Perlmutter2,Spergel1,Komatsu} that our universe is now accelerating. As an accelerated expansion phase, we know from the Friedmann equation that the second order derivative of the scale factor $\ddot{a}>0$ or $\dot{a}$ must be an increasing function over the time evolution. That means the Hubble parameter $H(t)$ is a decreasing function over the evolution of cosmic time. So Fig.\ref{f2} tells us that $H(t)$ was approximately same in the early stage of the universe and then it follows gradually decreasing and finally keeps a constant nature during the late time acceleration of the universe. The decreasing behavior of the Hubble parameter with time may also escape the universe from phantom-like evolution. As the Hubble parameter is proportional to the energy density in the late times of the evolution, here we get the same natural behaviour of the Hubble parameter in the late times of the cosmic evolution.

On the other hand, one can observe that the evolution of deceleration parameter starts at $q=-1$ which represents the de-Sitter expansion phase and then it goes to deceleration phase through accelerating power law expansion phase. After that, it again goes to the de-Sitter expansion phase in the late time of cosmic evolution (Fig.\ref{f3}). So $q(t)$ is positive during decelerating expansion and negative during accelerating expansion of the universe. Note that, for eternally accelerated universe $q<0$, $q=0$ occurs in the transition line and for the super-accelerated expansion, $q<-1$. This kinds of behavior of the universe is also suggested by the observed cosmological evolution of the universe \cite{Dodelson}.

 




\section{Cosmological Models in $f(Q,T)$ Gravity}\label{V}
In this present section, we discuss various functional forms of $f(Q,T)$ gravity model to analyse the cosmological evolutions in the symmetric teleparallel equivalent general relativity. The choices of $f(Q,T)$ are as follows.



\subsection{Case~I~: $f(Q,T)=\a Q+\b T$.}
The first simple form of $f(Q,T)$ has 
\ben
f(Q,T)=\a Q+\b T,\label{21}
\een
where $\a$ and $\b$ are constants. So we obtain $F=F_Q=\a$, and $8\pi\bar{G}=f_T=\b$.
\\
Now using the definition of scale factor mentioned in \eqref{18} and \eqref{21} in \eqref{13}, \eqref{14} and \eqref{17}, we have the following expressions for the energy density $\r$, pressure $p$, and EOS parameter $\o$. 
  
\ben
&&\r=-\Big(\dfrac{1}{8\pi +2\b}\Big)\Big[3\a\big(c+dm\tanh(n-mt)\big)^2+\frac{\b d m^2\a}{8\pi +\b}~\text{sech}^2(n-mt)\Big],\label{22}
\\
&&p=\dfrac{3\a\big[c+dm\tanh(n-mt)\big]^2}{\big(8\pi +2\b\big)}-\dfrac{\big(16\pi +3\b\big)~dm^2\a~\text{sech}^2(n-mt)}{\big(8\pi +2\b\big)\big(8\pi +\b\big)},\label{23}
\\
&&\o=\dfrac{dm^2\big(16\pi+3\b\big)\text{sech}^2(n-mt)-3\big(8\pi+\b\big)\big[c+dm\tanh(n-mt)\big]^2}{3\big(8\pi+\b\big)\big[c+dm\tanh(n-mt)\big]^2+\b dm^2~\text{sech}^2(n-mt)}.\label{24}
\een
 
The plots of energy density $\r$, pressure $p$, and the equation-of-state parameter $\o$ with the cosmic time $t$ is shown in Figs.\ref{f4}, \ref{5}, and \ref{6}, respectively.


\begin{figure}[H]
  \centering
  \includegraphics[width=8.5 cm]{rhovst_1.eps}
  \caption{Plot of energy density $(\r)$ as a function of cosmic time $(t)$ for $c=0.97,d=1,m=0.735,$ \& $n=10$ with $\a=0.1$ and $\b=-59.1$.}
  \label{f4}
\end{figure}

\begin{figure}[H]
  \centering
  \includegraphics[width=8.5 cm]{pvst_1.eps}
  \caption{Plot of pressure $(p)$ as a function of cosmic time $(t)$ for $c=0.97,d=1,m=0.735,$ \& $n=10$ with $\a=0.1$ and $\b=-59.1$.}
  \label{f5}
\end{figure}

\begin{figure}[H]
  \centering
  \includegraphics[width=8.5 cm]{omegavst_1.eps}
  \caption{Plot of EOS parameter $(\o)$ as a function of cosmic time $(t)$ for $c=0.97,d=1,m=0.735,$ \& $n=10$ with $\a=0.1$ and $\b=-59.1$.}
  \label{f6}
\end{figure}




\subsection{Case~II~: $f(Q,T)=u Q^{\e+1}+v T$.}
As another example of a cosmological model in the $f(Q,T)$ gravity, we consider a generic function, given by
\ben
f(Q,T)=u Q^{\e+1}+v T,\label{25}
\een
where $u,~v$ and $\e$ are constants. In this case, $F=f_Q=(\e+1)u Q^\e=(\e+1)u~6^\e H^{2\e}$ and $8\pi\bar{G}=v$. The generalised energy density $\r$, pressure $p$ and EOS parameter $\o$ in this case are given by

  
\ben
&&\r=-\Big(\dfrac{1}{8\pi +2v}\Big)\big[c+dm\tanh(n-mt)\big]^{2\e}(1+2\e)u 6^\e\Bigg[3\big[c+dm\tanh(n-mt)\big]^2+\frac{v(1+\e)}{(8\pi +v)}~dm^2~\text{sech}^2(n-mt)\Bigg],\non\\
\label{26}
\\
&&p=\dfrac{3u~6^\e(1+2\e)\big[c+dm\tanh(n-mt)\big]^2}{\big(8\pi +2v\big)}-\dfrac{6^\e (16\pi+3v)(1+\e)(1+2\e)~u dm^2~\text{sech}^2(n-mt)}{\big(8\pi +2v\big)\big(8\pi +v\big)}\Big[c+dm\tanh(n-mt)\Big]^{2\e},\non\\
\label{27}
\\
&&\o=-\dfrac{(8\pi+v)\big[c+dm\tanh(n-mt)\big]^{-2\e}3\big[c+dm\tanh(n-mt)\big]^2}{dm^2v(1+\e)~\text{sech}^2(n-mt)+3(8\pi+v)\big[c+dm\tanh(n-mt)\big]^2}\non\\
&&+\dfrac{\big[c+dm\tanh(n-mt)\big]^{-2\e}dm^2(16\pi+3v)(1+\e)~\text{sech}^2(n-mt)\big[c+dm\tanh(n-mt)\big]^{2\e}}{dm^2v(1+\e)~\text{sech}^2(n-mt)+3(8\pi+v)\big[c+dm\tanh(n-mt)\big]^2}.\label{28}
\een
 
It should be mentioned that for $\e=0$, the model in \eqref{25} goes to the first example of $f(Q,T)$ model. All the corresponding energy density, pressure and EOS parameter in \eqref{26}, \eqref{27} and \eqref{28} reduces to the \eqref{22}, \eqref{23} and \eqref{24}, respectively. So the choice of $f(Q,T)=u Q^{\e+1}+v T$ behaves as a more general one. The profiles of energy density, pressure and the EOS parameter vs. cosmic time with taking $\e=0.45$ is represented in Fig.\ref{f7}, \ref{f8} and \ref{f9}, respectively.


\begin{figure}[H]
  \centering
  \includegraphics[width=8.5 cm]{rhovst_2.eps}
  \caption{Plot of energy density $(\r)$ as a function of cosmic time $(t)$ for $c=0.97,d=1,m=0.735,$ \& $n=10$ with $u=0.1$ and $v=-59.1$.}
  \label{f7}
\end{figure}

\begin{figure}[H]
  \centering
  \includegraphics[width=8.5 cm]{pvst_2.eps}
  \caption{Plot of pressure $(p)$ as a function of cosmic time $(t)$ for $c=0.97,d=1,m=0.735,$ \& $n=10$ with $u=0.1$ and $v=-59.1$.}
  \label{f8}
\end{figure}

\begin{figure}[H]
  \centering
  \includegraphics[width=8.5 cm]{omegavst_2.eps}
  \caption{Plot of EOS parameter $(\o)$ as a function of cosmic time $(t)$ for $c=0.97,d=1,m=0.735,$ \& $n=10$ with $u=0.1$ and $v=-59.1$.}
  \label{f9}
\end{figure}

Now from the energy density profiles in the two cases, it is observed that the energy density is high in the early time of the universe and then gradually decreases and set to zero. On the other hand, the pressure $p$ lies in the negative range to suffice the acceleration of the universe which could satisfy the dark energy dominated accelerating universe. For the behaviour of equation of state parameter $\o$, it is clear that $\o$ takes the value exactly $-1$ for the first case (Fig.\ref{f6}), whereas it takes value close to $-1(\approx -0.66)$ for the second example (Fig.\ref{f9}) in the very early time, after that it smoothly raises to its maximum value and then starts decreases again. From Fig.\ref{f6}, it is observed that $\o$ takes it's maximum value which is nearly close to $1/3$ and then it converges to exactly $-1$ to produce the recent observation \cite{Riess1,Perlmutter1,Perlmutter2,Spergel1,Komatsu}. The same case is happening for the general case in Fig.\ref{f9} under the consideration of the constant $\e=0.45$, but here $\o$ decreases to more negative values $\sim -3.8$ to serve the huge late-times acceleration which could be generated by the negative pressure. Thus, we model a complete cosmic evolutionary history of the universe having inflation, radiation, matter and dark energy dominated eras in order by using a simple ansatz choice of scale factor. Obviously the ranges of the plot profiles of energy density, pressure and the equation of state parameter for the different cases are the different. This motivate to work with the different kinds of model consideration in $f(Q,T)$ gravity to study the cosmological scenarios. Finally from Figs.\ref{f6} and \ref{f9}, it is also worthy to mentioned that the universe starts with the acceleration smoothly, goes to the deceleration phase, and finally returns to its late phase of accelerated expansion, produced by dark energy.






\section{Various Energy Conditions}\label{VI}
The energy conditions are the expressions of the linear relationships between the energy density $\r$ and pressure $p$. It has a very crucial significance to understand the singularity theorem, to describe the nature of geodesics, to study the properties of black holes. The energy conditions are also the essential tools to study the geodesics of the Universe from the cosmological points of view. The physical conditions on the energy-momentum tensor are basically the energy conditions, which is obtained by making a connection between the Raychaudhuri equation \cite{Ray1,Ray2,Ray3} and the Einstein field equation.\\
From Raychaudhuri equation and by the constraints of gravity to be attractive in nature, we have the following two energy conditions, namely strong energy condition (SEC) and null energy condition (NEC) with a time-like vector field $v^\m$ and a null vector $k^\m$, respectively, i.e.,
\ben
&&\big(T_{\m\n}-\frac{1}{2}T g_{\m\n}\big)v^\m v^\n\geq 0\implies \r+3p\geq 0~,~~~\textbf{SEC}\label{29}
\\
&& T_{\m\n}k^\m k^\n\geq 0\implies \r+p\geq 0~,~~~\textbf{NEC}\label{30}
\een
On the other hand the weak energy condition (WEC) can be achieved by replacing null vector field $k^\a$ by a time-like vector $v^\a$ and the dominant energy condition (DEC) says that the matter should follow the time-like or null world lines.
\ben
&&\r\geq 0~~with~~\r+p\geq 0~,~~~\textbf{WEC}\label{31}
\\
&&\r\geq 0~~with~~\r\pm p\geq 0~,~~~\textbf{DEC}\label{32}
\een
where $T_{\m\n}$ is the energy-momentum tensor and $g_{\m\n}$ is the fundamental metric tensor ; $T=T^{\m\n}g_{\m\n}$, being the trace. Energy conditions in $f(Q,T)$ gravity have been studied in \cite{Simran1,Simran2}, but here we proceed to construct all the energy conditions on the background of \eqref{18}.



\subsection{Case-A : $f(Q,T)=\a Q+\b T$}
For the simplest choice of the $f(Q,T)=\a Q +\b T$ gravity model (where $\a,~\b$ are constants), we already have the expressions for the energy density $\r$ and pressure $p$ (\eqref{22}, \eqref{23}). As we know the energy density must be a positive quantity, here we chose the model parameters $\a=0.1,~\b=-59.1$ to ensure the energy density positive and the equation of state parameter as per the observations. Due to the importance of the energy conditions, we observed the acceptable ranges of model parameters $\a$ and $\b$, where the variation of $\a$ is from $9.1$ to $29.1$ and for $\b$, it is from $-25.55$ to $-55.55$.

Among all the energy conditions, it is easy to see from Fig.\ref{f10} that the strong energy condition (SEC) is violating on the cosmological scale, which is obvious according to the recent data observations of the accelerating Universe \cite{Barcelo, Moraes1}. The variation in $\a$ and $\b$ results in the variation of SEC behavior. Since the EOS parameter $\o$ is negative means $\r+3p<0$ and thus $\ddot{a}>0$, which implies there is a violation of SEC at present.



  

\begin{figure}[H]
     \centering
     \begin{subfigure}[b]{0.45\textwidth}
         \centering
         \includegraphics[width=0.96\textwidth]{sec_alpha_1.eps}
         \caption{Variation of $\a$ with $t$}
         \label{10a}
     \end{subfigure}
     \hfill
     \begin{subfigure}[b]{0.45\textwidth}
         \centering
         \includegraphics[width=0.96\textwidth]{sec_beta_1.eps}
         \caption{Variation of $\b$ with $t$}
         \label{10b}
     \end{subfigure}
\caption{Variation of SEC with $\a$ and $\b$}
        \label{f10}
\end{figure}


\begin{figure}[H]
     \centering
     \begin{subfigure}[b]{0.45\textwidth}
         \centering
         \includegraphics[width=0.96\textwidth]{nec_alpha_1.eps}
         \caption{Variation of $\a$ with $t$}
         \label{11a}
     \end{subfigure}
     \hfill
     \begin{subfigure}[b]{0.45\textwidth}
         \centering
         \includegraphics[width=0.96\textwidth]{nec_beta_1.eps}
         \caption{Variation of $\b$ with $t$}
         \label{11b}
     \end{subfigure}
\caption{Variation of NEC with $\a$ and $\b$}
        \label{f11}
\end{figure}

 On the other hand, the NEC and DEC both are obeying in this model (Figs.\ref{f11} and \ref{f12}). We have also observed the behavior of NEC, which is a partial condition of WEC, as we have shown the behavior of energy density in Fig.\ref{f4}. So it is obvious that the validation of NEC and energy density together results in the validation of WEC.

\begin{figure}[H]
     \centering
     \begin{subfigure}[b]{0.45\textwidth}
         \centering
         \includegraphics[width=0.96\textwidth]{dec_alpha_1.eps}
         \caption{Variation of $\a$ with $t$}
         \label{12a}
     \end{subfigure}
     \hfill
     \begin{subfigure}[b]{0.45\textwidth}
         \centering
         \includegraphics[width=0.96\textwidth]{dec_beta_1.eps}
         \caption{Variation of $\b$ with $t$}
         \label{12b}
     \end{subfigure}
\caption{Variation of DEC with $\a$ and $\b$}
        \label{f12}
\end{figure}

The complete story of all the energy conditions is depicted together in Fig.\ref{f13}. From Fig.\ref{f13}, it is observed that the NEC and WEC satisfy for the present model but SEC violates. Since we've already discussed the profile of $\o$ previously, it starts with acceleration, then smoothly goes to the deceleration phase, and again it returns to the accelerated phase. The exciting thing is getting the same results for SEC, NEC and DEC in Figs.\ref{f10}, \ref{f11} and \ref{f12}.

  

\begin{figure}[H]
  \centering
  \includegraphics[width=11cm]{ecs_all_macro_1.eps}
  \caption{Plot of Energy Conditions as a function of cosmic time $t$ for $c=0.97$, $d=1$, $m=0.735$, $n=10$, $\a=0.1$, \& $\b=-59.1$}
  \label{f13}
\end{figure}

 

\subsection{Case-B : $f(Q,T)=uQ^{\e+1}+v T$}
For the generic choice of $f(Q,T)=u Q^{\e+1} +v T$ gravity model (where $u,~v$ are arbitrary constants), we already have the expressions for the energy density $\r$ and pressure $p$ (\eqref{26}, \eqref{27}). Like the previous case, similarly we chose the model parameters $u=0.1,~v=-59.1$ with $\e=0.45$. Due to the importance of the energy conditions, again we observed the acceptable ranges of model parameters $u$ is from $9.1$ to $29.1$ and for $v$, it is from $-25.55$ to $-55.55$.
\\
Again among all the energy conditions, the strong energy condition is violated on cosmological scale (Fig.\ref{f14}). The variation in $u$ and $v$ results in the variation of SEC behavior. Since the EOS parameter $\o$ is negative means $\r+3p<0$, so that there is a violation of SEC at present epoch. The negative behavior of SEC executes the accelerated expansion of the universe. 

\begin{figure}[H]
     \centering
     \begin{subfigure}[b]{0.45\textwidth}
         \centering
         \includegraphics[width=0.96\textwidth]{sec_u_2.eps}
         \caption{Variation of $u$ with time $t$}
         \label{14a}
     \end{subfigure}
     \hfill
     \begin{subfigure}[b]{0.45\textwidth}
         \centering
         \includegraphics[width=0.96\textwidth]{sec_v_2.eps}
         \caption{Variation of $v$ with time $t$}
         \label{14b}
     \end{subfigure}
\caption{Variation of SEC with model parameters $u,~v$ and cosmic time $t$}
        \label{f14}
\end{figure}

\begin{figure}[H]
     \centering
     \begin{subfigure}[b]{0.45\textwidth}
         \centering
         \includegraphics[width=0.96\textwidth]{nec_u_2.eps}
         \caption{Variation of $u$ with $t$}
         \label{15a}
     \end{subfigure}
     \hfill
     \begin{subfigure}[b]{0.45\textwidth}
         \centering
         \includegraphics[width=0.96\textwidth]{nec_v_2.eps}
         \caption{Variation of $v$ with $t$}
         \label{15b}
     \end{subfigure}
\caption{Variation of NEC with model parameters $u,~v$ and cosmic time $t$}
        \label{f15}
\end{figure}

\begin{figure}[H]
     \centering
     \begin{subfigure}[b]{0.45\textwidth}
         \centering
         \includegraphics[width=0.96\textwidth]{dec_u_2.eps}
         \caption{Variation of $u$ with $t$}
         \label{16a}
     \end{subfigure}
     \hfill
     \begin{subfigure}[b]{0.45\textwidth}
         \centering
         \includegraphics[width=0.96\textwidth]{dec_v_2.eps}
         \caption{Variation of $v$ with $t$}
         \label{16b}
     \end{subfigure}
\caption{Variation of DEC with model parameters $u,~v$ and cosmic time $t$}
        \label{f16}
\end{figure}

On the other hand, the NEC and DEC, both do not violate in this model. So the behaviour of NEC and DEC is always positive (Figs.\ref{f15} and \ref{f16}). Therefore, the validation of NEC and energy density together results in the validation of WEC.  

Similarly, for the choice of $f(Q,T)=uQ^{1.45}+vT$ gravity model with the constants $u=0.1$ and $v=-59.1$, the portrait of all the energy conditions is shown in Fig.\ref{f17}. The same result is happening here i.e., the violation of SEC with the non-violation of NEC and DEC, which satisfy the recent observation data of the accelerated expansion.

\begin{figure}[H]
  \centering
  \includegraphics[width=11cm]{ecs_all_macro_2.eps}
  \caption{Plot of Energy Conditions as a function of cosmic time $t$ for $c=0.97$, $d=1$, $m=0.735$, $n=10$, $u=0.1$, \& $v=-59.1$}
  \label{f17}
\end{figure}






\section{Scalar Field Description}\label{VII}
In this section, we try to construct a scalar field description for the choice of scale factor \eqref{18} in $f(Q,T)$ gravity model for two different examples. This description is a more fundamental formulation for our study. To explain the early acceleration of the universe i.e., the inflationary framework, the most appropriate and popular way is to consider a scalar field, which is known as inflaton, managed by a specific potential. The action for the dynamics of inflation can be written as
\ben
\mathcal{A}=\int d^4 x \sqrt{-g}\Big[\frac{R}{2\kappa}+L^{matt}_{\phi}\Big].
\label{33}
\een
Here $g$ is the determinant of the metric tensor $g_{\m\n}$, $R$ is the Ricci scalar, $\kappa=8\pi$ is a constant ($G=1$). The matter Lagrangian $L^{matt}_{\phi}$ is containing an expression, which is minimally coupled of the inflaton field $\phi(t)$ to gravity, defined as follows.
\ben
L^{matt}_{\phi}=-\frac{1}{2}g^{\m\n}\p_\m \phi\p_\n \phi -V(\phi),
\label{34}
\een
where $V(\phi)$ is the potential of the scalar field $\phi(t)$ that can depends on one or more free parameters \cite{Lyth}.

The energy density $\r_{\phi}$ and the pressure $p_{\phi}$ in the scalar field model description can be described as
\ben
\r_{\phi}=\frac{1}{2}{\dot{\phi}}^2(t)+V(\phi),~~~~~~~p_{\phi}=\frac{1}{2}{\dot{\phi}}^2(t)-V(\phi),\label{35}
\een
where one can assume that the energy–momentum tensor of the inflaton field may behaves like a perfect fluid having a linear barotropic equation of state as $p_{\phi}=w\r_{\phi}$. On the other hand, the trace of the energy-momentum tensor in this scalar field description reads
\ben
T_{\phi}\equiv -\r_{\phi}+3p_{\phi}={\dot{\phi}}^2(t)-4V(\phi).\label{36}
\een



\subsection{Case-A}
For the linear model $f(Q,T)=\a Q+\b T$ using \eqref{35}, \eqref{15} and \eqref{16} can be rewritten as
\ben
&&3H^2=-\dfrac{(8\pi+\b){\dot{\phi_{1}^2}}(t)+2V(\phi_1)(8\pi+2\b)}{2\a}\label{37}
\\
&&4\dot{H}+3H^2=\dfrac{3(8\pi+\b){\dot{\phi_{1}^2}}(t)-2V(\phi_1)(8\pi+2\b)}{2\a}\label{38}
\een
From \eqref{37} and \eqref{38}, one can simply find out the relation between the Hubble parameter and the inflaton field $\phi_1$ as
\ben
\dot{H}=\frac{{\dot{\phi}}^2_{1}}{2\a}(8\pi+\b),\label{39}
\een
and we exclude the case for $\b=-8\pi$. Again, using \eqref{37} and \eqref{39}, we obtain the modified Klein-Gordon equation in this scalar field model as
\ben
\ddot{\phi_1}+3H\dot{\phi_1}+\frac{(8\pi+2\b)}{(8\pi+\b)}\frac{dV}{d\phi_1}=0.\label{40}
\een
Now by solving \eqref{39} with $\a=0.1$ and $\b=-59.1$, one may have the following.
\ben
&&\phi_1(t)=-0.0767\tan^{-1}\big[\sinh(n-mt)\big],\label{41}
\\
&&H(\phi_1)=c-dm\sin(13.0378~\phi_{1}),\label{42}
\\
&&V(\phi_1)=0.0005\Big[6\big\{c-dm\sin(13.0378~\phi_1)\big\}^2-2~dm^2\cos^2(13.0378~\phi_1)\Big].\label{43}
\een

\begin{figure}[H]
  \centering
  \includegraphics[width=8.5 cm]{v_phi1.eps}
  \caption{Plot of $V(\phi_1)$ with $\phi_1$ for $c=0.97$, $d=1$, $m=0.735$, \& $n=10$}
  \label{f18}
\end{figure}

This potential $V(\phi_1)$ may be seen as a combination of sine-cosine function. The parametric plot of $V(\phi_1)$ vs. $\phi_1$ has been given in Fig.\ref{f18}. It is clear from Fig.\ref{f18} that during the early times, $V(\phi_1)$ is slowly varying with the inflaton field $\phi_1$ and at the end of the inflation, the field is slowly rolling down and it becomes constant in the dark energy dominated era i.e., in the late-times. We may conclude the behaviour of constant potential implies the cosmological constant-like behavior during the late-times of the universe's evolution.



\subsection{Case-B}
Following the same procedure as in Case-A, we obtain the effective energy density, effective pressure in more generic $f(Q,T)=u Q^{\e+1}+v T$ as follows.
\ben
&&3H^2\equiv 8\pi\r_{eff}=\dfrac{-(8\pi+v){\dot{\phi_{2}^2}}(t)-2V(\phi_2)(8\pi+2v)}{2u~(1+2\e)~(8\pi)^{\e}~\big({\dot{\phi_{2}^2}}+2V\big)^{\e}},\label{44}
\\
&&2\dot{H}+3H^2\equiv -8\pi p_{eff}=\dfrac{(8\pi+v){\dot{\phi_{2}^2}}(t)-2V(\phi_2)(8\pi+2v)-\e\big((8\pi+v){\dot{\phi_{2}^2}}+2V(\phi_2)(8\pi+2v)\big)}{2u~(1+\e)(1+2\e)~(8\pi)^{\e}~\big({\dot{\phi_{2}^2}}+2V\big)^{\e}},\label{45}
\een
and from \eqref{44}, \eqref{45}, one can simply obtain an important relation i.e.,
\ben
H^{2\e}\dot{H}=\dfrac{(8\pi+v){\dot{\phi}}^2_{2}}{2u~6^{\e}(1+2\e)(1+\e)}.\label{46}
\een
By considering a binomial expansion and neglecting the higher order terms in \eqref{46} and excluding $8\pi=-v$ case, we've solved \eqref{46} numerically with $u=0.1,~v=-59.1$ and $\e=0.45$. The expressions for the scalar field in terms of cosmic time $t$ is given by,
\ben
\phi_{2}(t)=0.1874\Big[0.3421\text{sech}(n-mt)-\tan^{-1}\big[\sinh(n-mt)\big]\Big].\label{47}
\een
Due to that extra term in the expression of $\phi_{2}(t)$, it is difficult to express the Hubble parameter in terms of $\phi_{2}$ only. But to proceed further, we introduced the scalar field $\phi_{1}(t)$ from the first model in our second model and did a comparison study between both of these models. So by using \eqref{41}, \eqref{47} can be written as,
\ben
\phi_{2}(t)=\Big[0.0641\text{sech}(n-mt)+ 2.4432 \phi_{1}\Big].\label{48}
\een
In this case, the Hubble parameter and the potential can also be written as,
\ben
&&H=c+dm\sqrt{1-{\Big(\frac{\phi_{2}-2.4432\phi_{1}}{0.0641}\Big)}^2},\label{49}
\\
&&V=-\dfrac{u(1+2\e)6^{\e}H^{2\e}}{(2+\e)(8\pi+2v)}\Big[(1+\e)(2\dot{H}+3H^2)+3H^2+(1+\e)\e\dot{H}\Big].\label{50}
\een

From the scalar field solution for our second model, one can observe that an additional term arises due to the change in the Lagrangian $f(Q,T)$ compared to the first model.  After some numerical analysis, it is observed that our model-1 is able to describe the inflationary cosmology better than model-2 due to its' linearity in the nature of Lagrangian. Further, we have seen that the scalar field solution for the first model dominates the second model because the scalar field description for the second model depends on the first model. In addition, we found that the potential is symmetric in nature for the second model, i.e., it is increasing initially and attends a maximum value when the inflation ends and later converges towards zero. In contrast, we may say that our linear model is the best compared to the non-linear model. 






\section{Cosmic Evolution in $\omega-\omega^{\prime}$ Phase Plane}\label{VIII}
Cosmological evolution can also be presented in the well-known $\o-\o^{\prime}$ phase space representation. In this representation, by taking the scale factor as a parameter along the dynamics of the evolutionary path, one may define,
\ben
\o^{\prime}=\frac{d\o}{d\ln a}.\label{51}
\een
The evolutionary representation in this phase plane may be characterized into two regions, namely, thawing region and freezing region. The criterion for the identification of the thawing region is $\o^{\prime}>0$ and $\o<0$, and for the freezing region, it is given by $\o^{\prime}<0$ with $\o<0$. In the thawing region, $\o$ is an increasing function of time, so that the evolution of the universe terminate into de-Sitter like stage. On the other hand, $\o$ decreases on time evolution in the freezing model. So it's behaviour is asymptotic and it will be depended on the shape of the potential \cite{Caldwell,Linder,Sch}. Fig.\ref{f19} represents the evolution behaviour of the universe in $\o-\o^{\prime}$ phase space. In this Fig.\ref{f19}, the blue curve represents our model one i.e., $f(Q,T)=\a Q+\b T$, where the red trajectory define for the choice of the second model, $f(Q,T)=u Q^{\e+1}+v T$ with $\e=0.45$.

From Fig.\ref{f19} (for the blue trajectory), it is clear that the universe starts with the acceleration smoothly in the thawing region and then produces its second phase of accelerated expansion in the freezing region during late times. In between these two regions, the EOS parameter shows some positive behaviour (goes with the maximum value of $\o\sim \frac{1}{3}$) in the decelerating era, where $\o^{\prime}$ will be first increasing and then decreasing. Similarly same story is happening for the red curve also, when we have considered the contribution of $Q$ in the model is $1.45$. So the freezing region is larger compared to the thawing region due to the higher non-metricity effect in the model. As a whole, one may say that during early times, the universe exhibits thawing behavior, and now the universe is producing freezing behavior.

\begin{figure}[H]
  \centering
  \includegraphics[width=11cm]{phase.eps}
  \caption{(color online) The evolutionary tracks in $\o-\o^{\prime}$ phase space for $c=0.97$, $d=1$, $m=0.735$, $n=10$, $\a\equiv u=0.1$, \& $\b\equiv v=-59.1$}
  \label{f19}
\end{figure}








\section{Statefinder Diagnostics}\label{IX}
As our universe is ever accelerating with an unknown things, called dark energy, the unknown nature of dark energy arises many problems in modern cosmology. So there always exists a question that what is dark energy? Many dark energy models such as $\Lambda$CDM, HDE, SCDM, CG, Quintessence, k-essence have been proposed to understand the most intriguing nature of dark energy. And these dark energy models have different behaviours in comparison to each other. The $\left\lbrace r,s\right\rbrace$ parametrization technique is used to distinguish all these kind of dark energy models \cite{Sahni,Alam}.\\
Define $r$ and $s$ as follows.
\ben
&&r=\frac{\dot{\ddot{a}}}{aH^3},\label{52}
\\
&&s=\frac{2}{3}\frac{r-1}{2q-1},\label{53}
\een
where $q\neq \frac{1}{2}$. Now the different pairs of $\left\lbrace r,s\right\rbrace$ represents different kinds of dark energy models.

\begin{itemize}
\item[(a).] For the $\Lambda$CDM model, the pair is $\left\lbrace r=1,s=0\right\rbrace$.

\item[(b).] For the HDE model, the pair is $\left\lbrace r=1,s=\frac{2}{3}\right\rbrace$..

\item[(c).] For the SCDM model, the pair is $\left\lbrace r=1,s=1\right\rbrace$. 

\item[(d).] For the Quintessence model, the pair is $\left\lbrace r<1,s>0\right\rbrace$.

\item[(e).] For the CG model, the pair is $\left\lbrace r>1,s<0\right\rbrace$. 
\end{itemize}

We need to study the convergence and divergence nature of the trajectory of the $r-s$ curve corresponding to any cosmological dark energy models. So clearly the deviation from $\left\lbrace 1,0 \right\rbrace$ represents the deviation from $\Lambda$CDM model. Furthermore from the observation \cite{Albert,Albert2}, the values of $r$ and $s$ could be concluded. Obviously it is worthy of describing the various dark energy models in the near future.

Now using \eqref{18} in \eqref{52}, \eqref{53} we can rewrite the $r,s$ as,

\ben
&&r=\frac{c^3+d m \tanh (n-m t) \left\lbrace c^2+(d+1) m \tanh (n-m t) [3 c+(d+2) m \tanh (n-m t)]-(3 d+2) m^2\right\rbrace-3 c d m^2}{[c+d m \tanh (n-m t)]^3},\non
\\
\label{54}
\\
&&s=\frac{2 d m^2 \text{sech}^2(n-m t) [3 c+(3 d+2) m \tanh (n-m t)]}{3 [c+d m \tanh (n-m t)] \left\lbrace c^2+d m \tanh (n-m t) [6 c+(3 d+2) m \tanh (n-m t)]-2 d m^2\right\rbrace}.\non
\\
\label{55}
\een

\begin{figure}[H]
  \centering
  \includegraphics[width=9.5 cm]{rs.eps}
  \caption{$r-s$ parametric plot for $c=0.97,d=1,m=0.735,$ \& $n=10$.}
  \label{f20}
\end{figure}

\begin{figure}[H]
  \centering
  \includegraphics[width=6.5 cm]{rq.eps}
  \caption{$r-q$ parametric plot for $c=0.97,d=1,m=0.735,$ \& $n=10$.}
  \label{f21}
\end{figure}

In Fig.\ref{f20}, the parametrization of $r$ and $s$ shown in $(r,s)$ plane and the arrow mark represents the direction of the trajectory. From Fig.\ref{f20}, it is clearly observed that, initially, the trajectory diverges from $\Lambda$CDM model, and later, it converges to $\Lambda$CDM model. Also, the evolution of the trajectory completely lies in the Quintessence era. On the other hand, we have shown the parametrization of $r$ and $q$ in Fig.\ref{f21}. From Fig.\ref{f21}, we observed that our model starts with the de-Sitter universe, and initially, it goes to Chaplygin gas, represented by $r>1$ and after that, it comes to Quintessence and finally again converges to the de-Sitter universe.






\section{Conclusion}\label{X}
In this manuscript, we have aimed to explore a novel cosmological scenario in the context of $f(Q,T)$ gravity. As we know, our universe is passed through various evolutionary phases from early to the present day, and there are many cosmological models have been presented to discuss those scenarios separately in the gravitational theories. But still, we are in search of a cosmological model that can explain the complete evolution process of the universe from early inflation to late-time cosmic acceleration. In this regard, we have presented two $f(Q,T)$ cosmological models and discussed their matter evolution profiles and various energy conditions.

The equation of state parameter ($\omega$) plays a significant role in describing the various properties of the fluid in a cosmological model. Because different values of ($\omega$) present different matter profiles such as for radiation-dominated ($\omega=1/3$), matter-dominated ($\omega=0$), quintessence ($1/3<\omega<-1$), $\Lambda$CDM ($\omega=-1$), and phantom ($\omega<-1$). Therefore, we have examined the equation of state parameter over the cosmic time evolution for both models. It is observed that both models have shown accelerated expansion during the early-time and late-time evolution process. We also observed that model-I shows the $\Lambda$CDM type expansion for early-time and late-times, whereas model-II shows quintessence type for early-time and phantom-like for late-times. The sufficient amount of negative pressure of the universe in this model may be responsible for this accelerated expansion during the late-times. The complete profiles of the EOS parameter suggested that the evolution process of the universe started with early-time inflation and then went to its’ decelerated phase. Later, it shows the second phase of the accelerated expansion of the universe.

Furthermore, we have presented various energy conditions to check the viability of our cosmological models. From all the energy condition profiles, we have seen that the energy densities are positive throughout the evolution process, and WEC, NEC, and DEC are satisfied for both models. Also, we have seen that SEC is violated during the early and late-times evolution, which is an agreement for accelerated expansion. But, both models satisfied SEC during the decelerated phase of the universe. In addition, we have re-examined the dark energy behavior of our cosmological solution through $\o-\o^{\prime}$ phase space plane which is characterized by thawing and freezing regions and also by statefinder diagnostics. This behavior may also be confirmed with the scalar field description for the models in the background of $f(Q,T)$ gravity. We also observed that the potential $V(\phi_{1})$ is having constant value with varying $\phi_{1}$, which indicates the cosmological-constant-like behavior in the late times of cosmic expansion. Besides this, we also observed that our first model sounds good for an explanation of inflationary epochs than the second one due to the high non-linear effects in the considered Lagrangian, presents in the second model. 

In conclusion, our models have presented a complete evolution process of the universe from early to late-times in the context of $f(Q,T)$ gravity. We hope this study will shed some light on the idea of describing the whole evolution process of the universe through a cosmological model. In the future, it would be interesting to see the observational test of these types of models, and also, it may discover some new physics. 
 





\section{Acknowledgement}
S.D is thankful to Raj Kumar Das (senior research fellow, physics and applied mathematics unit, Indian Statistical Institute, Kolkata) for encouraging the aspects of the symmetric teleparallel theory of gravity.







\begin{thebibliography}{0}

\bibitem{Riess1}
A.G. Riess et al., Astron. J. \textbf{116}, 1009 (1998).

\bibitem{Perlmutter1}
S. Perlmutter et al., Astrophys. J. \textbf{517}, 565 (1999).

%\bibitem{de}
%P. deBernardis et al., Nature 404, 955 (2000).

%\bibitem{Colless}
%M. Colless et al., Mon. Not. R. Astron. Soc. 328, 1039 (2001).

\bibitem{Perlmutter2}
S. Perlmutter et al., Astrophys. J. \textbf{598}, 102 (2003).

\bibitem{Spergel1}
D.N. Spergel et al., Astrophys. J. Suppl. Ser. \textbf{148}, 175 (2003).

%\bibitem{Peiris}
%H.V. Peiris et al., Astrophys. J. Suppl. Ser. 148, 213 (2003).

%\bibitem{Tegmark}
%M. Tegmark et al., Phys. Rev. D 69, 103501 (2004).

%\bibitem{Cole}
%S. Cole et al., Mon. Not. R. Astron. Soc. 362, 505 (2005).

%\bibitem{Springel}
%V. Springel et al., Nature 440, 1137 (2006).

%\bibitem{Astier}
%P. Astier et al., Astron. Astrophys. 447, 31 (2006).

%\bibitem{Riess2}
%A.G. Riess et al., Astrophys. J. 659, 98 (2007).

%\bibitem{Spergel2}
%D.N. Spergel et al., Astrophys. J. Suppl. Ser. 170, 377 (2007).

%\bibitem{Ade}
%P.A.R. Ade et al., Astron. Astrophys. 571, A16 (2014).

\bibitem{Komatsu}
E. Komatsu et al., Astrophys. J. Suppl. Ser. \textbf{180} 330 (2009).

\bibitem{Capoz}
S. Capozziello et al., Int. J. Modern Phys. D \textbf{28}, 1930016 (2019). 

\bibitem{Nojiri1}
S. Nojiri et al., Phys. Rev. D. \textbf{72}, 023003 (2005).

\bibitem{Nojiri2}
S. Nojiri et al., Phys. Lett. B \textbf{686} (2010).

\bibitem{Bamba}
K. Bamba et al., Eur. Phys. J. C \textbf{76}, 18 (2016).

\bibitem{Astas}
A.V. Astashenok et al., Phys. Lett. B \textbf{713} (2012). 

\bibitem{Ratra}
B. Ratra et al., Phys. Rev. D \textbf{37}, 3406 (1988).

\bibitem{Caldwell1}
R.R.Caldwell et al., Phys. Rev. Lett. \textbf{80}, 1582 (1988).

\bibitem{Buchert}
T. Buchert, Gen. Relativ. Gravit. \textbf{32}, 105 (2000).

\bibitem{Picon}
C.A. Picon et al., Phys. Rev. D \textbf{63}, 103510 (2001).

\bibitem{Tomita}
K. Tomita et al., Mon. Not. R. Astron. Soc. \textbf{326}, 287 (2001).

\bibitem{Milton}
K.A. Milton et al., Gravit. Cosmol. \textbf{9}, 66 (2003).

\bibitem{Hunt}
P. Hunt et al., Mon. Not. R. Astron. Soc. \textbf{401}, 547 (2010).

\bibitem{Easson}
D. Easson et al., Phys. Lett. B \textbf{696}, 273 (2011).

\bibitem{Radicella}
N. Radicella et al., Gen. Relativ. Gravit. \textbf{44}, 685 (2012).

\bibitem{Pavn}
D. Pavón et al., Gen. Relativ. Gravit. \textbf{45}, 63 (2013).

\bibitem{Pandey}
B. Pandey, Mon. Not. R. Astron. Soc. \textbf{471}, L77 (2017).

%\bibitem{Pandey2}
%B. Pandey, Mon. Not. R. Astron. Soc. 485, L73 (2019).

\bibitem{Palia}
A. Paliathanasis et al., Phys. Rev. D \textbf{94}, 023525 (2016).

\bibitem{Soti}
T.P. Sotiriou et al., Rev. Mod. Phys. \textbf{82}, 451 (2010).

\bibitem{Har}
T. Harko et al., Phys. Rev. D \textbf{84}, 024020 (2011).


\bibitem{Bengochea}
G.R. Bengochea et al., Phys. Rev. D \textbf{79}, 124019 (2009).

\bibitem{Ferraro}
R. Ferraro et al., Phys. Rev. D \textbf{75}, 084031 (2007).

\bibitem{Linder1}
E.V. Linder, Phys. Rev. D \textbf{81}, 127301 (2010).

\bibitem{Cai}
Y.F. Cai et al., Rep. Progr. Phys. \textbf{79}, 106901 (2016).

\bibitem{Nester}
J.M. Nester et al., Chin. J. Phys. \textbf{37}, 113 (1999).

\bibitem{Lavinia}
J.B. Jimenez et al., Phys. Rev. D \textbf{98}, 044048 (2018).

\bibitem{Nisha}
N. Godani et al., Int. J. Geom. Methods Mod. Phys. \textbf{09}, 2150134 (2021).

\bibitem{Xu}
Y. Xu et al., Eur. Phys. J. C \textbf{80}, 449 (2020).

\bibitem{Gad}
G.N. Gadbail et al., Eur. Phys. J. C \textbf{81}, 1088 (2021).

\bibitem{Simran3}
S. Arora et al., Eur. Phys. J. C \textbf{81}, 555 (2021).

\bibitem{Antonio}
A. Nájera et al., J. Cosmol. Astropart. Phys. \textbf{03}, 020 (2022).

\bibitem{Maryam}
M. Shiravand et al., Phys. Dark Univ. \textbf{37} 101106, (2022).

\bibitem{Koussour}
M. Koussour et al., Int. J. Geom. Methods Mod. Phys. \textbf{16}, 2250115 (2022).

 
\bibitem{Lalke}
A.R. Lalke et al., Int. J. Geom. Methods Mod. Phys \textbf{09}, 2350131 (2023).

\bibitem{Simran1}
S. Arora et al., Phys. Scr. \textbf{95}, 095003 (2020).

\bibitem{Simran2}
S. Arora et al., Phys. Dark Univ. \textbf{31}, 100790 (2021).

\bibitem{Tayde}
M. Tayde et al., Chin. Phys. C \textbf{46} 115101 (2022). 

\bibitem{Asutosh}
S. Mandal et al., Int. J. Geom. Methods Mod. Phys. \textbf{05}, 2350084 (2023).

\bibitem{gaurav} G.N. Gaidbail et al., Phys. Lett. B \textbf{838} 137710 (2023).
 
\bibitem{Weyl/1918}
H.H. Weyl et al., Phys. Math. Kl. \textbf{465}, 1 (1918).

\bibitem{Di}
P.A.M. Dirac, Proc. R. Soc. A \textbf{333}, 403 (1973).


\bibitem{Von}
H.H.V. Borzeszkowski et al., Gen. Relativ. Gravit. \textbf{29}, 455 (1997).

\bibitem{D}
D. Puetzfeld, Class. Quantum Gravity \textbf{19}, 3263 (2002).

\bibitem{Z}
Z. Haghani et al., Class. Quantum Gravity \textbf{32}, 215016 (2015).

\bibitem{M1}
M. Novello et al., Phys. Rep. \textbf{463}, 127 (2008).

\bibitem{T}
T. Ortin, Gravity and Strings, Cambridge Monographs on Mathematical Physics (Cambridge University Press, Cambridge, England, 2015).

\bibitem{L}
L. Järv et al., Phys. Rev. D \textbf{97}, 124025 (2018).

\bibitem{J}
J.W. Maluf et al., Ann. Phys. \textbf{525}, 339 (2013).

\bibitem{Z1}
Z. Haghani et al., J. Cosmol. Astropart. Phys. \textbf{10}, 061 (2012).

\bibitem{Z2}
Z. Haghani et al., Phys. Rev. D \textbf{88}, 044024 (2013).

\bibitem{Adak}
M. Adak et al., Int. J. Mod. Phys. A \textbf{28}, 1350167 (2013).

\bibitem{Mol}
I. Mol et al., Advances in Applied Clifford Algebras \textbf{27}, 2607 (2017).

\bibitem{Conroy}
A. Conroy et al.,  Eur. Phys. J. C \textbf{78}, 923 (2018). 

\bibitem{Kovisto}
T. Koivisto et al., Int. J. Geo. Meth. Mod. Phys. \textbf{15}, 1840006 (2018).

\bibitem{Beltran}
J.B. Jimenez et al., J. Cosmol. Astropart. Phys. \textbf{08}, 039 (2018).

\bibitem{Soudi}
I. Soudi et al., Phys. Rev. D \textbf{100}, 044008 (2019).

\bibitem{Dia}
K.F. Dialektopoulos et al., Eur. Phys. J. C \textbf{79}, 606 (2019).

\bibitem{Lu}
J. Lu et al., Eur. Phys. J. C \textbf{79}, 530 (2019).

\bibitem{Laz} 
R. Lazkoz et al., Phys. Rev. D \textbf{100}, 104027 (2019).

\bibitem{Lavi2}
J.B. Jiménez et al., Phys. Rev. D \textbf{101}, 103507 (2020).

\bibitem{Sanjay1} 
S. Mandal et al., Phys. Rev. D \textbf{102}, 024057 (2020).

\bibitem{Sanjay2} 
S. Mandal et al., Phys. Rev. D \textbf{102}, 124029 (2020).

\bibitem{Sneha}
O. Sokoliuk et al., Eur. Phys. J. Plus \textbf{137}, 1077 (2022). 

\bibitem{Mustafa1}
G. Mustafa et al., Phys. Lett. B \textbf{821}, 136612 (2021).

\bibitem{Hassan1}
Z. Hassan et al., Fortschr. Phys. \textbf{69}, 2100023 (2021).


\bibitem{Raja}
R. Solanki et al., Phys. Dark Univ. \textbf{32}, 100820 (2021).

\bibitem{De}
A. De et al., Eur. Phys. J. C \textbf{82}, 72 (2022).

\bibitem{Raja2}
R. Solanki et al., Phys. Dark Univ. \textbf{36}, 100996 (2022).

\bibitem{Sanjay3}
S. Mandal et al., Phys. Lett. B \textbf{823}, 136786 (2021).

\bibitem{Sanjay4}
S. Mandal et al., Phys. Dark Univ. \textbf{35}, 100934, (2022).

\bibitem{Harko}
T. Harko et al., Phys. Rev. D \textbf{98}, 084043 (2018).

\bibitem{Moraes}
P.H.R.S. Moraes et al., Eur. Phys. J. C \textbf{76}, 60 (2016).

\bibitem{Dodelson}
S. Dodelson, Modern Cosmology, Academic Press, Amsterdam, (2003).

\bibitem{Ray1}
A. Raychaudhuri, Phys. Rev. \textbf{98}, 1123 (1955).

\bibitem{Ray2}
A. Raychaudhuri, Z. Astrophys. \textbf{43}, 161 (1957).

\bibitem{Ray3}
A. K. Raychaudhuri, S. Banerji, A. Banerjee, General Relativity,
Astrophysics and Cosmology (Astronomy \& Astrophysics Library, Springer, 2003).

\bibitem{Barcelo}
C. Barcelo et a., Int. J. Mod. Phys. D \textbf{11}, 1553 (2002).

\bibitem{Moraes1}
P.H.R.S. Moraes et al., Eur. Phys. J. C \textbf{77}, 480 (2017).

\bibitem{Lyth}
D.H. Lyth et al., Phys. Rep. \textbf{314}, 1 (1999).

\bibitem{Caldwell}
R.R. Caldwell et al., Phys. Rev. Lett. \textbf{95}, 141301 (2005). 

\bibitem{Linder}
E.V. Linder, Phys. Rev. D \textbf{73}, 063010 (2006). 

\bibitem{Sch}
R.J. Scherrer et al., Phys. Rev. D \textbf{77}, 083515 (2008).

\bibitem{Sahni}
V. Sahni et al., JETP Lett. \textbf{77}, 201 (2003).

\bibitem{Alam}
U. Alam et al., Mon. Not. R. Astron. Soc. \textbf{344}, 1057 (2003).

\bibitem{Albert}
J. Albert et al., SNAP Collaboration, arXiv:astro-ph/0507460.

\bibitem{Albert2}
J. Albert et al., SNAP Collaboration, arXiv:astro-ph/0507459.


\end{thebibliography}




\end{document}


